\par\noindent\textbf{(1)} 定义 \(\gamma_{A,j}=c_j+c_j^\dagger\)、\(\gamma_{B,j}=i(c_j^\dagger-c_j)\)，它们满足 \(\gamma^\dagger=\gamma\)、\(\{\gamma_{a,j},\gamma_{b,k}\}=2\delta_{ab}\delta_{jk}\)。反解 \(c_j=(\gamma_{A,j}+i\gamma_{B,j})/2\)，故每个复费米模由两个实Majorana分量构成。
\par\noindent\textbf{(2)} 在 \(\mu=0,\Delta=t\) 的标准Kitaev链中，Hamilton量可整理为 \(H=it\sum_{j=1}^{N-1}\gamma_{B,j}\gamma_{A,j+1}\)（整体正负随Hamilton规范）。只有相邻位点Majorana相互耦合，而最左 \(\gamma_{A,1}\) 与最右 \(\gamma_{B,N}\) 不进入Hamilton量。
\par\noindent\textbf{(3)} 令 \(f=(\gamma_{A,1}+i\gamma_{B,N})/2\)，则 \(\{f,f^\dagger\}=1\)。因两端Majorana不与 \(H\) 耦合，\([H,f]=0\)，两个态 \(|0_f\rangle\) 与 \(f^\dagger|0_f\rangle\) 能量相同但总费米奇偶不同，构成拓扑边缘零模的两重简并。
\par\noindent\textbf{(4)} 当 \(|\mu|<2t\) 且 \(\Delta\ne0\) 时理想无限链处于拓扑相，其边缘Majorana具有指数局域尾；有限长链的两端尾部重叠产生 \(\Delta E\propto e^{-N/\xi}\) 的小劈裂。超过 \(|\mu|=2t\) 的体能隙闭合转变后，零模不再受相同拓扑保护。
